\section*{Chapter \ref{ch:m1:futures-contracts-and-exchanges} --- Futures Contracts and Their Exchanges}

\begin{solution}{exo:m1:futures-contracts-and-exchanges:1}
$(6\,012.50 - 5\,987.25)/0.25 = 101$ \omterm{def:m1:futures-contracts-and-exchanges:tick}{ticks}; $101 \times \$12.50 \times 12 =
\$15\,150$ lost.
\end{solution}

\begin{solution}{exo:m1:futures-contracts-and-exchanges:2}
$111 + 24.5/32 = 111.765625$ and $112 + 8/32 = 112.25$: a rise of $0.484375$
point, 31 \omterm{def:m1:futures-contracts-and-exchanges:tick}{ticks} of $1/64$, worth $31 \times \$15.625 = \$484.375$.
\end{solution}

\begin{solution}{exo:m1:futures-contracts-and-exchanges:3}
EURO STOXX 50 future, March (H) of the first year ending in 7 that is not
past: March 2027. Last trading day: the third Friday, 19 March 2027, at 12:00
CET. \texttt{ESH6} cannot be March 2026, which has expired: read today it
means March 2036, which is not listed. A symbol with a one-digit year is
ambiguous without a date, which is why reference data carries full expiry
dates.
\end{solution}

\begin{solution}{exo:m1:futures-contracts-and-exchanges:4}
\omterm{def:m1:what-a-trading-firm-does:market-maker}{Market maker}: the minimum spread, one \omterm{def:m1:futures-contracts-and-exchanges:tick}{tick}, is 1.85 basis points against
0.42, so each round trip captured earns more, and the queue at each price is
long: priority in the queue, hence speed and the matching algorithm, decides
who earns it. Prices move less often, so quotes are safer for longer. Fund:
\euro50 million is 926 contracts; crossing a \omterm{def:m1:futures-contracts-and-exchanges:tick}{one-tick} spread costs half a
\omterm{def:m1:futures-contracts-and-exchanges:tick}{tick}, $926 \times \text{\euro}5 = \text{\euro}4\,630$, 0.93 basis point, where the
same trade in a contract with the finer \omterm{def:m1:futures-contracts-and-exchanges:tick}{tick} would cost about a quarter of
that.
\end{solution}

\begin{solution}{exo:m1:futures-contracts-and-exchanges:5}
$300\,000/21\,000 = 14.3$ times; a move of 7\% loses the deposit. The deposit
is the minimum the clearing house demands today. Losses must be paid in cash
every day as \omterm{def:m1:clearing-and-settlement:margin}{variation margin}, and the \omterm{def:m1:clearing-and-settlement:margin}{initial margin} itself is raised when
volatility rises: the money needed to \emph{hold} the position through a bad
week is several times the deposit (\cref{ch:m1:margin}).
\end{solution}

\begin{solution}{exo:m1:futures-contracts-and-exchanges:6}
After A buys 5 from B: A $+5$, B $-5$, \omterm{def:m1:futures-contracts-and-exchanges:oi}{open interest} 5. C buys 3 from A: A
$+2$, C $+3$, B $-5$; A has passed 3 contracts to C, \omterm{def:m1:futures-contracts-and-exchanges:oi}{open interest} 5. B buys
2 from C: B $-3$, C $+1$; both closed, \omterm{def:m1:futures-contracts-and-exchanges:oi}{open interest} 3. D buys 4 from A: A
$-2$, D $+4$; A closes 2 and opens 2 short, D opens 4: \omterm{def:m1:futures-contracts-and-exchanges:oi}{open interest} 5
(longs: C 1, D 4; shorts: A 2, B 3). Volume 14, \omterm{def:m1:futures-contracts-and-exchanges:oi}{open interest} 5.
\end{solution}

\begin{solution}{exo:m1:futures-contracts-and-exchanges:7}
Largest residual: 14.29\% with the \$300\,000 contract (at \$1.05 million,
half a contract unhedged) and 1.00\% with the small one. The residual is at
most half a contract: $150\,000/V$ and $15\,000/V$, that is $14.3\%$ and
$1.4\%$ at $V = \$1.05$ million; the data respect both bounds.
\end{solution}

\begin{solution}{exo:m1:futures-contracts-and-exchanges:8}
Every open contract has a buyer and a seller: buyers never outnumber sellers
in contracts. A rise in \omterm{def:m1:futures-contracts-and-exchanges:oi}{open interest} says that new longs \emph{and} new
shorts opened positions; the price rose because the buyers were the more
impatient side, which the price already says. What the data can tell:
whether a move comes with new positions (rising \omterm{def:m1:futures-contracts-and-exchanges:oi}{open interest}) or with
liquidation (falling); where the roll stands; and, combined with the
regulator's weekly report, which class of trader holds which side.
\end{solution}

\begin{solution}{pb:m1:futures-contracts-and-exchanges:1}
\textbf{1.} \$300\,000 and \euro54\,000.
\textbf{2.} $1.15 \times 187\,000\,000/300\,000 = 716.83$: sell 717. Residual
$215.05 - 215.10 = -\$50\,000$: over-hedged by a sixth of a contract.
\textbf{3.} $0.90 \times 43\,000\,000/54\,000 = 716.67$: sell 717; residual
$-\text{\euro}18\,000$.
\textbf{4.} December 2026, \texttt{ESZ6} and \texttt{FESXZ6}, last trading
day Friday 18 December 2026, exactly three months away. The E-mini settles
against the \emph{opening} quotation of that day and the FESX at noon: the
hedge ends some hours before the close of the last day.
\textbf{5.} $717 \times 300\,000 \times 5\% = \$10\,755\,000$.
\textbf{6.} Futures: $717 \times 300\,000 \times 2\% = +\$4\,302\,000$. Sleeve:
$1.15 \times 2\% \times 187$ million $= -\$4\,301\,000$.
\textbf{7.} The futures gain is paid in cash the next morning; the loss on
the shares is unrealised. On a good day it is the reverse: the fund must
\emph{pay} \$4.3 million of \omterm{def:m1:clearing-and-settlement:margin}{variation margin} in cash against an unrealised
gain, and must hold or borrow that cash.
\textbf{8.} $4\,302\,000 - 2.9\% \times 187\,000\,000 = -\$1\,121\,000$: the
sleeve fell 0.6\% more than its beta predicted. This is the idiosyncratic and
factor risk of the portfolio, which an index hedge does not remove.
\textbf{9.} $0.10 \times 187$ million $= \$18.7$ million of index exposure,
374 times the rounding residual.
\textbf{10.} $717 \times 0.05 \times \$50 = \$1\,792.50$.
\textbf{11.} The currency: the sleeve is a euro asset for a dollar fund, and
the FESX gains and losses are in euros too. The hedge removes the equity
beta, not the \euro43 million of euro exposure (in fact it adds a small
euro exposure on its margin and P\&L).
\textbf{12.} A small noise with zero mean: the two prices are taken at the
same moment by different mechanisms, a thirty-second average for the future
and a single auction price for each share. On ordinary days it is a basis
point or two; on index-rebalance and expiry days, when the closing auction
moves, it is larger, and it reverses the next day.
\textbf{13.} Liquidity and cost: the E-mini is the cheapest hedge per dollar.
It gives up the match: a small-capitalisation sleeve hedged with a
large-capitalisation index keeps the size spread, which can move several
percent in a quarter. A small-capitalisation index future hedges better and
costs more to trade.
\textbf{14.} A short future \emph{earns} the implied financing rate less
dividends: the fund receives the benchmark rate plus the futures' \omterm{def:m1:delta-one-instruments:swap}{funding
spread} on the hedged notional, in exchange for giving up the equity
premium.
\textbf{15.} Stock selection (intended), beta error, size and sector
mismatch, currency on the European sleeve, liquidity for \omterm{def:m1:clearing-and-settlement:margin}{variation margin},
and the hours each day when futures trade and shares do not.
\textbf{16.} The future cannot trade below the limit while the stock market
is halted; when both reopen the next limit applies. The hedge still marks
correctly at settlement, but cannot be adjusted during the halt, and the
shares cannot be sold either.
\textbf{17.} Barely: the residual would fall from \$50\,000 to at most
\$15\,000, on a position whose beta uncertainty is worth \$18.7 million. It
matters for portfolios a hundred times smaller.
\textbf{18.} Cost and reversibility: selling \$187 million of shares costs
tens of basis points, realises taxes, and abandons the stock selection; 717
futures cost about a basis point and leave the portfolio intact.
\textbf{19.} \textbf{Sell 717 E-minis; residual $-\$50\,000$.}
\textbf{20.} It standardises everything except the price, the quantity and
the choice of which contract approximates the risk one actually has.
\end{solution}

\begin{solution}{iq:m1:futures-contracts-and-exchanges:1}
The multiplier is \$50, the \omterm{def:m1:futures-contracts-and-exchanges:tick}{tick} 0.25 index point, so a \omterm{def:m1:futures-contracts-and-exchanges:tick}{tick} is \$12.50. At
an index of 6\,000 the contract is worth \$300\,000; a 1\% move is \$3\,000 per
contract.
\iqlookfor{instant recall, and the habit of converting to notional.}
\end{solution}

\begin{solution}{iq:m1:futures-contracts-and-exchanges:2}
Volume counts contracts traded during a period; \omterm{def:m1:futures-contracts-and-exchanges:oi}{open interest} counts
contracts in existence at a moment. A trade between two new participants
adds to both; a trade in which one side closes leaves \omterm{def:m1:futures-contracts-and-exchanges:oi}{open interest}
unchanged; a trade between two closers reduces it. A day trader generates
volume and no \omterm{def:m1:futures-contracts-and-exchanges:oi}{open interest}.
\iqlookfor{the three cases.}
\end{solution}

\begin{solution}{iq:m1:futures-contracts-and-exchanges:3}
As integer \omterm{def:m1:futures-contracts-and-exchanges:tick}{ticks}, with the \omterm{def:m1:futures-contracts-and-exchanges:tick}{tick} size and multiplier held once, exactly, in
reference data. Comparisons and keys are exact, arithmetic is fast, prices
off the grid cannot exist, and money is obtained by one multiplication at
the edge. Decimal strings such as 0.01 or prices in thirty-seconds are
parsed once on input. \omterm{def:m1:futures-contracts-and-exchanges:tick}{Tick} sizes change (and differ for spreads and
options), so the conversion belongs to the instrument and the date.
\iqlookfor{integers, one conversion point, and awareness that \omterm{def:m1:futures-contracts-and-exchanges:tick}{tick} sizes
change.}
\end{solution}

\begin{solution}{iq:m1:futures-contracts-and-exchanges:4}
A share is issued by a company and any venue may trade it; a future is
created by an exchange, which owns the contract and clears it in its own
clearing house. \omterm{def:m1:futures-contracts-and-exchanges:oi}{Open interest} is not fungible across clearing houses, so a
competitor's identical contract starts with no liquidity and no margin
offsets against the incumbent's other products. Liquidity then attracts
liquidity.
\iqlookfor{clearing as the lock-in, not the trading.}
\end{solution}

\begin{solution}{iq:m1:futures-contracts-and-exchanges:5}
Find out how the series was built: on which day it switches contract and
whether it is adjusted. Compute each day's return from two prices of the
\emph{same} contract, switching contracts on a roll date chosen by a rule
(volume or \omterm{def:m1:futures-contracts-and-exchanges:oi}{open interest} crossing, or fixed days before expiry), so that the
roll gap never enters a return. Decide between ratio- and
difference-adjusted levels according to use; check settlement against last
trade; align the settlement time with any other series used with it.
\iqlookfor{returns from same-contract prices; an explicit roll rule.}
\end{solution}

\begin{solution}{iq:m1:futures-contracts-and-exchanges:6}
Which index: the one closest to the portfolio, traded off against liquidity.
Beta: estimated how, over what window, with what error. Contracts: $\beta
V/(mF)$, rounded. Expiry: \omterm{def:m1:futures-contracts-and-exchanges:code}{front month} and roll, or the month that covers the
horizon. Execution: a schedule against the hedge's urgency, perhaps at the
close to match the portfolio's valuation. Margin: the deposit, plus cash for
\omterm{def:m1:clearing-and-settlement:margin}{variation margin} on rallies. Currency, if any. Monitoring: beta drift,
residual factor exposures, the roll. Exit: the same in reverse.
\iqlookfor{the cash for \omterm{def:m1:clearing-and-settlement:margin}{variation margin}, which most candidates forget.}
\end{solution}
